%%%%%%%%%%%%%%%%%%%%%%%%%%%%%%%%%%%%%%%%%%%%%
%%% ! WARNUNG ! %%%%%%%%%%%%%%%%%%%%%%%%%%%%%
%%%%%%%%%%%%%%%%%%%%%%%%%%%%%%%%%%%%%%%%%%%%%
%%%  Diese Datei bitte nur bearbeiten,    %%%
%%%   wenn du ein LaTeX-Experte bist      %%%
%%%             U N D                     %%%
%%%  die Vorlage unbedingt ändern willst  %%%
%%%%%%%%%%%%%%%%%%%%%%%%%%%%%%%%%%%%%%%%%%%%%
%%%%%%%%%%%%%%%%%%%%%%%%%%%%%%%%%%%%%%%%%%%%%
%
%%%%%%%%%%%%%%%%%%%%%%%%%%%%%%%%%%%%%%%%%%%%%%%%%%%%%%%%%%%%%%%%%%%%%%%%%%%%%%%%
%%% DATEI-INFO %%%%%%%%%%%%%%%%%%%%%%%%%%%%%%%%%%%%%%%%%%%%%%%%%%%%%%%%%%%%%%%%%
%%%%%%%%%%%%%%%%%%%%%%%%%%%%%%%%%%%%%%%%%%%%%%%%%%%%%%%%%%%%%%%%%%%%%%%%%%%%%%%%
%%% Diese Datei richtet diverse Pakete und Befehle für die Vorlage ein %%%%%%%%%
%%%%%%%%%%%%%%%%%%%%%%%%%%%%%%%%%%%%%%%%%%%%%%%%%%%%%%%%%%%%%%%%%%%%%%%%%%%%%%%%
%
\def\inPreamble{}% Flag setzen: wird sind in der Präambel
%%%%%%%%%%%%%%%%%%%%%%%%%%%%%%%%%%%%%%%%%%%%%%%%%%%%%%%%%%%%%%%%%%%%%%%%%%%%%%%%
%%% Encoding %%%%%%%%%%%%%%%%%%%%%%%%%%%%%%%%%%%%%%%%%%%%%%%%%%%%%%%%%%%%%%%%%%%

\RequirePackage[utf8]{inputenc}
% UTF8 Codierte .tex-Dateien

\RequirePackage[T1]{fontenc}
% Benutzung der europäischen Schriftkodierung
% z.B. zur richtigen Darstellung von Umlauten


%%%%%%%%%%%%%%%%%%%%%%%%%%%%%%%%%%%%%%%%%%%%%%%%%%%%%%%%%%%%%%%%%%%%%%%%%%%%%%%%
%%% Werkzeuge %%%%%%%%%%%%%%%%%%%%%%%%%%%%%%%%%%%%%%%%%%%%%%%%%%%%%%%%%%%%%%%%%%

\RequirePackage{etoolbox}		% Einfacheres Programmieren in LaTeX
\RequirePackage{calculator}		% Rechnen (auch mit Längen)
\RequirePackage{xfp}			% Genaueres Rechnen (aber nicht mit Längen)
\RequirePackage{pbox}			% \pbox (Parbox flexibler Breite)


%%%%%%%%%%%%%%%%%%%%%%%%%%%%%%%%%%%%%%%%%%%%%%%%%%%%%%%%%%%%%%%%%%%%%%%%%%%%%%%%
%%% Variablen zur Anpassung durch den Benutzer %%%%%%%%%%%%%%%%%%%%%%%%%%%%%%%%%
%%%%%%%%%%%%%%%%%%%%%%%%%%%%%%%%%%%%%%%%%%%%%%%%%%%%%%%%%%%%%%%%%%%%%%%%%%%%%%%%
%%% >>> werden gesetzt in der Datei "PersoenlicheAngaben.tex" %%%%%%%%%%%%%%%%%%

\newbool{abbildungsverzeichnis}
\newbool{quellcodeverzeichnis}
\newbool{tabellenverzeichnis}
\newbool{symbolverzeichnis}
\newbool{akronymverzeichnis}
\newbool{abkuerzungsverzeichnis}
\newbool{glossar}

\newbool{verzeichnisseZusammenfassen}
\newbool{verzeichnisseImInhaltsverzeichnis}

\newbool{nichtZitiertInweiterfuehrendeLiteratur}
\newbool{weiterfuehrendeLiteratur}
\newbool{alleAutorenExplizitNennen}

\newbool{sperrvermerk}
\newbool{verlaengerung}
\newbool{danksagung}
\newbool{zusatzErklaerung}
\newbool{aufgabenstellung}
\newbool{eidesstattlicheVersicherung}
\newbool{kiErklaerung}

\newbool{color}
\newbool{linksMarkieren}

\newbool{doppelseitig}

\newlength{\bindekorrektur}
\setlength{\bindekorrektur}{1cm}

\newlength{\randAussen}
\setlength{\randAussen}{1.88cm} % 1.88 - 2.5cm (einseitig), 2.5 - 3.4 cm (doppelseitig)


%%%%%%%%%%%%%%%%%%%%%%%%%%%%%%%%%%%%%%%%%%%%%%%%%%%%%%%%%%%%%%%%%%%%%%%%%%%%%%%%
%%% Benutzer-Einstellungen anwenden %%%%%%%%%%%%%%%%%%%%%%%%%%%%%%%%%%%%%%%%%%%%

\input{Einstellungen}


%%%%%%%%%%%%%%%%%%%%%%%%%%%%%%%%%%%%%%%%%%%%%%%%%%%%%%%%%%%%%%%%%%%%%%%%%%%%%%%%
%%% Verzeichnis- & PDF-Verzeichnis-Einträge generieren %%%%%%%%%%%%%%%%%%%%%%%%%

%TODO Verwendung so anpassen, dass der Eintrag (=Linkziel) direkt beim entsprechenden Titel steht (=Titel auch mit \verzeichnisEintrag generieren und vom originalverzeichnis ausblenden?)
\newcommand{\verzeichnisEintrag}[2]{%
	% #1 Name
	% #2 Label (für PDF)
	\ifbool{verzeichnisseImInhaltsverzeichnis}{%
%		\ifbool{verzeichnisseZusammenfassen}{%
			%nichts
%		}{%
			\clearpage%
%		}%
		\phantomsection% damit die Seitenzahl auf jeden Fall stimmt
		\nopagebreak%
		\addcontentsline{toc}{chapter}{#1}% 
		\nopagebreak%
	}{%
		\pdfbookmark[0]{#1}{#2}% hier passte die Seitenzahl bisher immer, daher keine \phantomsection nötig
	}%
}

\newcommand{\verzeichnisEintragX}[3]{%
	% #1 Name
	% #2 Label (für PDF)
	\ifbool{verzeichnisseImInhaltsverzeichnis}{%
		%		\ifbool{verzeichnisseZusammenfassen}{%
		%nichts
		%		}{%
		\clearpage%
		%		}%
		\phantomsection% damit die Seitenzahl auf jeden Fall stimmt
		\nopagebreak%
		\addcontentsline{toc}{chapter}{\texorpdfstring{#1}{#2}}% 
		\nopagebreak%
	}{%
		\pdfbookmark[0]{#2}{#3}% hier passte die Seitenzahl bisher immer, daher keine \phantomsection nötig
	}%
}


%%%%%%%%%%%%%%%%%%%%%%%%%%%%%%%%%%%%%%%%%%%%%%%%%%%%%%%%%%%%%%%%%%%%%%%%%%%%%%%%
%%% Document-Class %%%%%%%%%%%%%%%%%%%%%%%%%%%%%%%%%%%%%%%%%%%%%%%%%%%%%%%%%%%%%

\ifbool{doppelseitig}{
	\def\onetwoside{twoside}
}{
	\def\onetwoside{oneside}
}

\documentclass[
	\onetwoside,% twoside: zweiseitig, oneside: einseitig
	titlepage, 	% Titelseite steht für sich allein
	12pt,		% Schriftgröße
	a4paper, 	% A4-Papier
]{report}



%%%%%%%%%%%%%%%%%%%%%%%%%%%%%%%%%%%%%%%%%%%%%%%%%%%%%%%%%%%%%%%%%%%%%%%%%%%%%%%%
%%% Sprachen %%%%%%%%%%%%%%%%%%%%%%%%%%%%%%%%%%%%%%%%%%%%%%%%%%%%%%%%%%%%%%%%%%%

\RequirePackage[%
	main=\hauptsprache, %Umsetzung der deutschen Sprache nach neuer Rechtschreibung
	\hauptsprache, % muss hier nochmal erwähnt werden, damit \autoref das auch mitbekommt
	\weitereSprachen %Sprachen, die nur Abschnittsweise (z.B. für Zitate) verwendet werden
	]{babel}

\RequirePackage{translations}

\DeclareTranslation{German}{Bibliography}{Literatur}
\DeclareTranslation{English}{Bibliography}{Bibliography}
\DeclareTranslationFallback{Bibliography}{Bibliography}

\DeclareTranslation{German}{FurtherReading}{Weiterführende Literatur}
\DeclareTranslation{English}{FurtherReading}{Further Reading}
\DeclareTranslationFallback{FurtherReading}{Further Reading}

\DeclareTranslation{German}{abbreviationsname}{Abkürzungen}
\DeclareTranslation{English}{abbreviationsname}{Abbreviations}
\DeclareTranslationFallback{abbreviationsname}{Abbreviations}

%%%%%%%%%%%%%%%%%%%%%%%%%%%%%%%%%%%%%%%%%%%%%%%%%%%%%%%%%%%%%%%%%%%%%%%%%%%%%%%%
%%% Schriftarten, Schriftsatz %%%%%%%%%%%%%%%%%%%%%%%%%%%%%%%%%%%%%%%%%%%%%%%%%%

\RequirePackage{lmodern} % Schriftart lmodern nutzen
\RequirePackage{microtype} % Besserer Schriftsatz

%%%%%%%%%%%%%%%%%%%%%%%%%%%%%%%%%%%%%%%%%%%%%%%%%%%%%%%%%%%%%%%%%%%%%%%%%%%%%%%%
%%% Zitate, Anführungszeichen %%%%%%%%%%%%%%%%%%%%%%%%%%%%%%%%%%%%%%%%%%%%%%%%%%

\RequirePackage[
	autostyle, % Zitierstil an die aktuelle Sprache anpassen
	german=quotes, % wenn Deutsch -> „Anführungszeichen“ benutzen 
]{csquotes} %wird auch von biblatex verwendet

\renewcommand{\mkcitation}[1]{\,#1} % Umdefiniert um doppelklammerung bei \textquote+\cite zu verhindern

%%%%%%%%%%%%%%%%%%%%%%%%%%%%%%%%%%%%%%%%%%%%%%%%%%%%%%%%%%%%%%%%%%%%%%%%%%%%%%%%
%%% Literaturverwaltung %%%%%%%%%%%%%%%%%%%%%%%%%%%%%%%%%%%%%%%%%%%%%%%%%%%%%%%%

\def\autorenAnzahlImLiteraturverzeichnis{4}
\ifbool{alleAutorenExplizitNennen}{
	\def\autorenAnzahlImLiteraturverzeichnis{99}
}{}

\RequirePackage[backend=bibtex8,  %Daten von Bibtex holen
			style=alphabetic, %In der Form [Aut22] referenzieren (Autorinitialen+Jahreszahl)
%			style=numeric, %In der Form [42] zitieren (Quellen durchnummerieren)
			sorting=nyt, % Sortierung: Name,Year,Title 
%			sorting=none, % Sortierung: Nach erstem Zitat
			sortcites=false, % Quellen im Zitat (um-)sortieren
			maxbibnames=\autorenAnzahlImLiteraturverzeichnis, %Anzahl maximal/minimal genannter Autoren im Literaturverzeichnis
			minbibnames=3, 
			maxcitenames=3, %Anzahl maximal/minimal genannter Autoren mit \citeauthor{...}
			mincitenames=3,
			maxsortnames=99, % Anzahl der Autoren, die in die Sortierung eingehen
			minsortnames=99,
			maxalphanames=4, % Maximale Anzahl Autoren, bei der jeder Autor mit einem Buchstaben in [ABCD] genannt wird
			minalphanames=3, % Anzahl Autoren, ab der [ABC+] verwendet wird
			]{biblatex}
			
% URLs im Literaturverzeichnis umbrechen, wenn nicht genug Platz vorhanden ist.
\setcounter{biburlnumpenalty}{9000} % An Ziffern
\setcounter{biburllcpenalty}{9000}  % An Kleinbuchstaben
\setcounter{biburlucpenalty}{9000}  % An Großbuchstaben

\DeclareNameAlias{default}{family-given}% Format: Nachname, Vorname
\DefineBibliographyStrings{ngerman}{ 
	andothers = {{et\,al\adddot}}, % et al.
}
\renewcommand*{\multinamedelim}{\addsemicolon\space} % Zeichen zwischen den Autoren
%\renewcommand*{\finalnamedelim}{\addsemicolon\space} % Zeichen zwischen den letzten beiden Autoren

\addbibresource{Verzeichnisse/Literatur.bib} % Literatur-Datei laden

\DeclareBibliographyCategory{cited} % registrieren, welche Quellen zitiert wurden
\AtEveryCitekey{%
	\addtocategory{cited}{\thefield{entrykey}}%
}

\AtEndDocument{% In manchen Installationen funktionierte das auch in der Präambel, aber so ist es sicherer
	\ifbool{nichtZitiertInweiterfuehrendeLiteratur}{%
		\nocite{*} % alle nicht zitierten Quellen mit aufnehmen
	}{%
		\explizitesNocite%
	}
}


%%%%%%%%%%%%%%%%%%%%%%%%%%%%%%%%%%%%%%%%%%%%%%%%%%%%%%%%%%%%%%%%%%%%%%%%%%%%%%%%
%%% Seitengröße festlegen %%%%%%%%%%%%%%%%%%%%%%%%%%%%%%%%%%%%%%%%%%%%%%%%%%%%%%

% Berechnung der seitlichen Seitenränder:
\def\randAussenCm{}
\def\bindekorrekturCm{}
\LENGTHDIVIDE{\randAussen}{1cm}{\randAussenCm} 			% auf cm normieren
\LENGTHDIVIDE{\bindekorrektur}{1cm}{\bindekorrekturCm} 	% auf cm normieren

\newlength{\randRechtsAussen}	% Als Länge anlegen (bessere Wiederverwendbarkeit)
\newlength{\randLinksInnen}

\setlength{\randRechtsAussen}{\fpeval{1*\randAussenCm} cm}
\setlength{\randLinksInnen}{\fpeval{\randAussenCm\ifbool{doppelseitig}{/2}{}+\bindekorrekturCm} cm} %TODO ggf. Formel ändern?

% Anwenden
\usepackage[
%	showframe, 				% Seitenränder anzeigen - zum Debuggen
	layout=a4paper,			% Papierformat
	headheight=1.2cm,		% durch Kopfzeile festgelegt
	top=3.55cm,				% Abstand Text <-> oberer Seitenrand
	left=\randLinksInnen,	% Abstand Text <-> linker Seitenrand
	right=\randRechtsAussen,% Abstand Text <-> rechter Seitenrand
	bottom=3.7cm 			% Abstand Text <-> unterer Seitenrand
]{geometry}


%%%%%%%%%%%%%%%%%%%%%%%%%%%%%%%%%%%%%%%%%%%%%%%%%%%%%%%%%%%%%%%%%%%%%%%%%%%%%%%%
%%% (R), (C), TM (vor allem für Zitate und das Literaturverzeichnis) %%%%%%%%%%%

\newcommand{\markRegistered}[1]{#1$^\text{\tiny{\textregistered\xspace}}$}
\newcommand{\markCopyrighted}[1]{#1$^\text{\tiny{\copyright\xspace}}$}
\newcommand{\markTrademark}[1]{#1$^\text{\tiny{\texttrademark\xspace}}$}
\usepackage[official]{eurosym}


\usepackage{bbding}%Häkchen für abgehakte Dinge
\newcommand{\ok}{\Checkmark}
\newcommand{\x}{\XSolidBrush}
\newcommand{\xg}{\textcolor{gray}{\x}}


%%%%%%%%%%%%%%%%%%%%%%%%%%%%%%%%%%%%%%%%%%%%%%%%%%%%%%%%%%%%%%%%%%%%%%%%%%%%%%%%
%%% Farben %%%%%%%%%%%%%%%%%%%%%%%%%%%%%%%%%%%%%%%%%%%%%%%%%%%%%%%%%%%%%%%%%%%%%

\usepackage[
	table, 		% Lade das colortbl package (Farben in Tabellen)
	showerrors,	% Nicht definierte Farbe verwendet -> Fehler erzeugen
]{xcolor}		% Farben
\selectcolormodel{\colormodel}% Konvertiere alle Farben in das angegebene Farbmodell

% TMDT-Farbpalette
\definecolor{TmdtDarkBlue}	{RGB}{1, 66, 98}    % #014262
\definecolor{TmdtLightBlue}	{RGB}{101, 190, 218} % #65BEDA
\definecolor{TmdtOrange}	{RGB}{255, 181, 99}  % #FFB563

\definecolor{pureK}			{cmyk}{0.00, 0.00, 0.00, 1.00} % K

\colorlet{AccentStrong}		{TmdtOrange}
\colorlet{AccentWeak}		{TmdtLightBlue!65!TmdtDarkBlue}

% Farbschema
\ifbool{color}{
	% wenn Farbe
	\colorlet{primaryTop}		{TmdtDarkBlue}
	\colorlet{primaryMiddle}	{TmdtDarkBlue}
	\colorlet{primaryBottom}	{TmdtLightBlue}
	\colorlet{secondaryTop}		{TmdtLightBlue}
	\colorlet{secondaryMiddle}	{TmdtLightBlue!65!TmdtDarkBlue}
	\colorlet{secondaryBottom}	{TmdtDarkBlue}
	\colorlet{tertiaryTop}		{TmdtOrange}
	\colorlet{tertiaryMiddle}	{TmdtOrange}
	\colorlet{tertiaryBottom}	{TmdtOrange}
	\colorlet{textBlack}		{black}
}{
	% wenn Graustufen
	\colorlet{primaryTop}		{pureK}
	\colorlet{primaryMiddle}	{pureK}
	\colorlet{primaryBottom}	{pureK}
	\colorlet{secondaryTop}		{pureK!70}
	\colorlet{secondaryMiddle}	{pureK!70}
	\colorlet{secondaryBottom}	{pureK!70}
	\colorlet{tertiaryTop}		{pureK!50}
	\colorlet{tertiaryMiddle}	{pureK!50}
	\colorlet{tertiaryBottom}	{pureK!50}
	\colorlet{textBlack}		{pureK}
}

% Kurzform für mittlere Farbe
\colorlet{primary}			{primaryMiddle}
\colorlet{secondary}		{secondaryMiddle}
\colorlet{tertiary}			{tertiaryMiddle}


% Spezialisierter Anwendungszweck
\colorlet{pagenumber}		{tertiaryBottom!50!textBlack}
\colorlet{headrule}			{primaryTop}
\colorlet{footnoteRule}		{secondary}
\colorlet{footrule}			{primaryBottom}

% TODO Formelnummern Style
\colorlet{captionLabel}		{pureK!80}
\colorlet{subcaptionLabel}	{pureK!70}

\colorlet{captionText}		{pureK!70}
\colorlet{subcaptionText}	{pureK!60}

\colorlet{listingBackground}{pureK!5}

%%%%%%%%%%%%%%%%%%%%%%%%%%%%%%%%%%%%%%%%%%%%%%%%%%%%%%%%%%%%%%%%%%%%%%%%%%%%%%%%
%%% Fußnotengestaltung %%%%%%%%%%%%%%%%%%%%%%%%%%%%%%%%%%%%%%%%%%%%%%%%%%%%%%%%%
%
\renewcommand{\footnoterule}{%
	\color{footnoteRule}{%
%	\textcolor{footnoteRule}{
% 	an der selben Stelle sorgt für Probleme mit umgebrochenen Listings.
%	Das liegt daran, dass \textcolor{}{} zusätzlich zu \color auch noch ein \leavevmode enthält.
%	Das \leavevmode wiederum bringt das Listing durcheinander, so dass z.B. die Zeilennummern direkt beim Umbruch 
%	jeweils vertauscht sind und oberhalb der untersten Zeile der vorderen Seite eine Lücke entsteht.
%	
		\vspace*{-3pt}%
		\hrule height 0.1em width 0.4\columnwidth%
		\vspace*{2.6pt}%
	}%
}%


%%%%%%%%%%%%%%%%%%%%%%%%%%%%%%%%%%%%%%%%%%%%%%%%%%%%%%%%%%%%%%%%%%%%%%%%%%%%%%%%
%%% Kopfzeilen und Fußzeilen, Seitenzahl %%%%%%%%%%%%%%%%%%%%%%%%%%%%%%%%%%%%%%%

\def\customFootrule{
	\renewcommand{\footrulewidth}{1.5mm}
	\renewcommand{\footrule}{%
		\begingroup
			\color{footrule}
			\hrule height \footrulewidth
			\vspace{0.5em}
		\endgroup
	}
}

\def\customHeadrule{
	\renewcommand{\headrulewidth}{1.5mm}
	\renewcommand{\headrule}{%
		\begingroup
			\color{headrule}
			\hrule height \headrulewidth 
		\endgroup
	}
}


\def\renewMarks{
	\renewcommand{\sectionmark}[1]{%
		\markright{%
			##1%~(\thesection)%
		}%
	}
	\renewcommand{\chaptermark}[1]{%
		\markboth{%
			##1% \chaptername~\thechapter%
		}{%
		}%
	}
}

\def\seitenNummer{%
	\color{pagenumber}\sffamily\bfseries\thepage{}%
}

\usepackage{fancyhdr} % Schönere Kopfzeilen
\def\vertikalZentrieren{%
	\vspace{-\footrulewidth}%
	\vspace{-1.275em}%
}

\ifbool{doppelseitig}{
	\def\rightodd{RO}
	\def\righteven{RE}
	\def\leftodd{LO}
	\def\lefteven{LE}
	\def\rightoddlefteven{RO,LE}
	\def\leftoddrighteven{LO,RE}
}{
	\def\rightodd{R}
	\def\righteven{R}
	\def\leftodd{R}
	\def\lefteven{R}
	\def\rightoddlefteven{R}
	\def\leftoddrighteven{L}
}

\newlength{\headerbuwlogowidth} % measured width of the BUW/TMDT logo block
\newlength{\headertextsecondaryheight} % measured height of the secondary text in its box
\newlength{\headerPrimaryWidth} % maximum width available for the Primary Header Text
\newlength{\headerSecondaryWidth}
\newbool{headerSecondaryTextFillsSpaceBeneathLogo}

\fancypagestyle{thesis-page-regular}{%
	\renewMarks%
	\fancyhf{}%clear
	\fancyhead[\rightoddlefteven]{%
		\buwandtmdtlogos{2em}%
	}%
	\fancyhead[\leftoddrighteven]{%
		\settowidth{\headerbuwlogowidth}{\buwandtmdtlogos{2em}}%
		\setlength{\headerSecondaryWidth}{\linewidth-\headerbuwlogowidth-1em}%
		%
		\settoheight{\headertextsecondaryheight}{\parbox[b]{\headerSecondaryWidth}{\sffamily\rightmark}}%
		\ifdimgreater{\headertextsecondaryheight}{1.9em}{%
			\setbool{headerSecondaryTextFillsSpaceBeneathLogo}{true}%
			\setlength{\headerPrimaryWidth}{\linewidth}%
		}{%
			\setbool{headerSecondaryTextFillsSpaceBeneathLogo}{false}%
			\setlength{\headerPrimaryWidth}{\headerSecondaryWidth}%
		}%
%		
		\parbox[b]{\headerPrimaryWidth}{% Für extrem lange Überschriften, damit die nicht ins Logo geraten
			\ifbool{doppelseitig}{% Rechts/Linksbündigkeit wiederherstellen
				\ifnumodd{\thepage}{\raggedright}{\raggedleft}%
			}{%
				\raggedright%
			}%
			\nouppercase{%
				\sffamily%
				\textbf{\leftmark}%
			}%
			\\\vspace{0.2em}%
			\parbox[b]{\headerSecondaryWidth}{% Für extrem lange Überschriften, damit die nicht ins Logo geraten
				\ifbool{doppelseitig}{% Rechts/Linksbündigkeit wiederherstellen
					\ifnumodd{\thepage}{\raggedright}{\raggedleft}%
				}{%
					\raggedright%
				}%
				\nouppercase{%
					\sffamily%
					\rightmark%
				}%
			}%
		}%
	}
	\fancyfoot[\rightodd]{\vertikalZentrieren\colorbox{white}{\seitenNummer}\hspace{1em}}
	\fancyfoot[\lefteven]{\vertikalZentrieren\hspace{1em}\colorbox{white}{\seitenNummer}}
	
	\customHeadrule
	\customFootrule
}

\fancypagestyle{plain}{%
	\renewMarks
	\fancyhf{} % clear
	\fancyfoot[\rightodd]{\vertikalZentrieren\colorbox{white}{\seitenNummer}\hspace{1em}}
	\fancyfoot[\lefteven]{\vertikalZentrieren\hspace{1em}\colorbox{white}{\seitenNummer}}
	
	\customHeadrule
	\customFootrule
}


%%%%%%%%%%%%%%%%%%%%%%%%%%%%%%%%%%%%%%%%%%%%%%%%%%%%%%%%%%%%%%%%%%%%%%%%%%%%%%%%
%%% Format der Überschriften %%%%%%%%%%%%%%%%%%%%%%%%%%%%%%%%%%%%%%%%%%%%%%%%%%%

\usepackage[%
	nobottomtitles*,	% Überschriften am unteren Seitenrand vermeiden
]{titlesec}
% 
\def\titleformatBase{\sffamily\bfseries}
\def\titleformatChapter{\Huge}
\def\titleformatSection{\LARGE}
\def\titleformatSubsection{\Large}
\def\titleformatSubsubsection{\large}
\def\titleformatParagraph{\normalsize}
\def\titleformatSubparagraph{\normalsize}

\def\titlepartChapter		{\titleformatChapter\ifbool{weAreInTheAppendix}{\Alph{chapter}}{\arabic{chapter}}}
\def\titlepartSection		{\ifnumcomp{\value{secnumdepth}}{>}{0}{\titleformatSection\arabic{section}}{}}
\def\titlepartSubsection	{\ifnumcomp{\value{secnumdepth}}{>}{1}{\titleformatSubsection\arabic{subsection}}{}}
\def\titlepartSubsubsection	{\ifnumcomp{\value{secnumdepth}}{>}{2}{\titleformatSubsubsection\arabic{subsubsection}}{}}
\def\titlepartParagraph		{\ifnumcomp{\value{secnumdepth}}{>}{3}{\titleformatParagraph\arabic{paragraph}}{}}
\def\titlepartSubparagraph	{\ifnumcomp{\value{secnumdepth}}{>}{4}{\titleformatSubparagraph\arabic{subparagraph}}{}}

% Platz zwischen Nummer und Titel
\def\titleSpaceSubparagraph	{{}}
\def\titleSpaceParagraph	{\hphantom{\titleformatParagraph\ifnumcomp{\value{secnumdepth}}{>}{4}{.}{}\titlepartSubparagraph}\titleSpaceSubparagraph}
\def\titleSpaceSubsubsection{\hphantom{\titleformatSubsubsection\ifnumcomp{\value{secnumdepth}}{>}{3}{.}{}\titlepartParagraph}\titleSpaceParagraph}
\def\titleSpaceSubsection	{\hphantom{\titleformatSubsection\ifnumcomp{\value{secnumdepth}}{>}{2}{.}{}\titlepartSubsubsection}\titleSpaceSubsubsection}
\def\titleSpaceSection		{\hphantom{\titleformatSection\ifnumcomp{\value{secnumdepth}}{>}{1}{.}{}\titlepartSubsection}\titleSpaceSubsection}
\def\titleSpaceChapter		{\hphantom{\titleformatChapter\ifnumcomp{\value{secnumdepth}}{>}{0}{.}{}\titlepartSection}\titleSpaceSection}

\newbool{weAreInTheAppendix}
\setbool{weAreInTheAppendix}{false}
\appto{\appendix}{\setbool{weAreInTheAppendix}{true}}
\def\titleNumberExtraDistance{.5cm}

\titleformat{\chapter} %command
	[hang] %shape
	{\titleformatBase} %format
	{%
		\titlepartChapter\titleSpaceChapter%
	} %label
	{\titleNumberExtraDistance} %sep
	{\titleformatChapter} %before-code
%	[after-code]
	
\titleformat{\section}
	[hang]
	{\titleformatBase}
	{%
		\color{textBlack!50}\titlepartChapter.%
		\color{textBlack}\titlepartSection%
		\titleSpaceSection%
	}
	{\titleNumberExtraDistance}
	{\titleformatSection}
	%[after-code]

\titleformat{\subsection}
	[hang]
	{\titleformatBase}
	{%
		\color{textBlack!33}\titlepartChapter.%
		\color{textBlack!50}\titlepartSection.%
		\color{textBlack}\titlepartSubsection%
		\titleSpaceSubsection%
	}
	{\titleNumberExtraDistance}
	{\titleformatSubsection}
	%[after-code]
	
\titleformat{\subsubsection}
	[hang]
	{\titleformatBase}
	{%
		\color{textBlack!20}\titlepartChapter.%
		\color{textBlack!33}\titlepartSection.%
		\color{textBlack!50}\titlepartSubsection.%
		\color{textBlack}\titlepartSubsubsection%
		\titleSpaceSubsubsection%
	}
	{\titleNumberExtraDistance}
	{\titleformatSubsubsection}
	%[after-code]
	
\titleformat{\paragraph}
	[runin]
	{\titleformatBase}
	{%
		\color{textBlack!15}\titlepartChapter.%
		\color{textBlack!20}\titlepartSection.%
		\color{textBlack!33}\titlepartSubsection.%
		\color{textBlack!50}\titlepartSubsubsection.%
		\color{textBlack}\titlepartParagraph%
		\titleSpaceParagraph%
	}
	{\titleNumberExtraDistance}
	{\titleformatParagraph}%
	%[after-code]
	
\titleformat{\subparagraph}
	[runin]
	{\titleformatBase}
	{%
		\color{textBlack!10}\titlepartChapter.%
		\color{textBlack!15}\titlepartSection.%
		\color{textBlack!20}\titlepartSubsection.%
		\color{textBlack!33}\titlepartSubsubsection.%
		\color{textBlack!50}\titlepartParagraph.%
		\color{textBlack}\titlepartSubparagraph%
		\titleSpaceSubparagraph}
	{\titleNumberExtraDistance}
	{\titleformatSubparagraph}%
	%[after-code]
	
%				command			left	before	after
\titlespacing*{\chapter}		{0pt}	{-2.1em}{1em}
\titlespacing*{\section}		{0pt}	{1em}	{1em}
\titlespacing*{\subsection}		{0em}	{1em}	{1em}
\titlespacing*{\subsubsection}	{0em}	{1em}	{1em}
\titlespacing*{\paragraph}		{0em}	{1em}	{1em}
\titlespacing*{\subparagraph}	{0em}	{1em}	{1em}

%%%%%%%%%%%%%%%%%%%%%%%%%%%%%%%%%%%%%%%%%%%%%%%%%%%%%%%%%%%%%%%%%%%%%%%%%%%%%%%%
%%% Grafiken %%%%%%%%%%%%%%%%%%%%%%%%%%%%%%%%%%%%%%%%%%%%%%%%%%%%%%%%%%%%%%%%%%%

\usepackage{graphicx}				% um Bilder einbinden zu können 
\usepackage{pdfpages}				% Ganzseitig externe PDFs einbinden (z.B. für Aufgabenstellung)

\usepackage{tikz}					% Zeichnen per Code


%%%%%%%%%%%%%%%%%%%%%%%%%%%%%%%%%%%%%%%%%%%%%%%%%%%%%%%%%%%%%%%%%%%%%%%%%%%%%%%%
%%% SI-Einheiten und Zahlendarstellung %%%%%%%%%%%%%%%%%%%%%%%%%%%%%%%%%%%%%%%%%

\usepackage[
	% Die Einstellungen für dieses Paket befinden sich in der Datei siunitx.cfg
]{siunitx}


%%%%%%%%%%%%%%%%%%%%%%%%%%%%%%%%%%%%%%%%%%%%%%%%%%%%%%%%%%%%%%%%%%%%%%%%%%%%%%%%
%%% Quellcode %%%%%%%%%%%%%%%%%%%%%%%%%%%%%%%%%%%%%%%%%%%%%%%%%%%%%%%%%%%%%%%%%%
\usepackage{listings} %Für Quelltexte

% manuell Übersetzen
\DeclareTranslation{German}{lstlistingname}{Quellcode}
\DeclareTranslation{English}{lstlistingname}{Sourcecode}
\DeclareTranslationFallback{lstlistingname}{\lstlistingname}
\renewcommand{\lstlistingname}{\GetTranslation{lstlistingname}}

\DeclareTranslation{German}{lstlistoflistingname}{Quellcodeverzeichnis}
\DeclareTranslation{English}{lstlistoflistingname}{List of Sourcecodes}
\DeclareTranslationFallback{lstlistoflistingname}{\lstlistlistingname}
\renewcommand{\lstlistlistingname}{\GetTranslation{lstlistoflistingname}}

\ifbool{color}{
	\def\listingBasicstyle		{\color{black}\ttfamily\small}
	\def\listingKeywordstyle	{\color{primary}\bfseries}
	\def\listingStringstyle		{\color{secondary}\slshape}
	\def\listingCommentstyle	{\color{AccentWeak}\fontseries{l}\selectfont}
	\def\listingEmphstyle		{\color{AccentStrong}\bfseries}
	\def\listingRulecolor		{\color{AccentWeak}}
	\def\listingNumbercolor		{\color{AccentWeak!60}}
}{
	\def\listingBasicstyle		{\color{pureK}\ttfamily\small}
	\def\listingKeywordstyle	{\color{pureK}\bfseries}
	\def\listingStringstyle		{\color{pureK}\slshape}
	\def\listingCommentstyle	{\color{pureK}\fontseries{l}\selectfont}
	\def\listingEmphstyle		{\color{pureK}\bfseries\underbar}
	\def\listingRulecolor		{\color{pureK!50}}
	\def\listingNumbercolor		{\color{pureK!26}}
}

\lstset{
	basicstyle	= \listingBasicstyle,		% Basisformat
	keywordstyle= \listingKeywordstyle,		% Schlüsselwörter
	stringstyle	= \listingStringstyle,		% Strings
	commentstyle= \listingCommentstyle,		% Kommentare
	emphstyle	= \listingEmphstyle,		% Hervorhebungen
%	
%	extendedchars=true,      	%Erweiterte Zeichensätze in Listings benutzen (aber nur Single-Byte!)   
	breaklines=true,            % Zeilen werden Umgebrochen
	breakatwhitespace=true,		% erlaube Zeilenumbrüche nur an Whitespace
	breakindent=0em,
	breakautoindent=true, % bei Zeilenumbruch automatisch einrücken
%
	prebreak=,
	postbreak=\mbox{$\hookrightarrow$\hspace{-0.34em}\hphantom{m}},
%	
	showstringspaces=true, 		% Leerzeichen in Strings zeigen
	showspaces=false,           % Leerzeichen anzeigen?
	showtabs=false,             % Tabs anzeigen?
	tabsize=3,					% Zeichenbreite eines Tabs
%	
	keepspaces=false,			% Erlaubt die Unterdrückung von Leerzeichen zugunsten eines besseren Spalten-Layouts
	columns=fixed,				% Zeichenbreite
%	
	captionpos=b, 				% Position der Beschriftung
	xleftmargin=2.2em, 			% Abstand zum Rand
	xrightmargin=0.35em,		%
	framexleftmargin=0em,		% Abstand im Rahmen
	framexrightmargin=0pt,		%
	framexbottommargin=2pt,		%
	framextopmargin=0pt,		%
	frame=tblr,					% Rahmen anzeigen
	frameround=tftf,			% t: abgerundete Ecke, f: eckige Ecke
	backgroundcolor=\color{listingBackground}, % Hintergrundfarbe
	rulecolor=\listingRulecolor,% Farbe der Umrandung
%
	numbers=left,
	numberstyle=\scriptsize\ttfamily\listingNumbercolor, %Stil der Zeilennummern
	stepnumber=1,
}

\lstset{literate= % Deutsche Umlaute und ß in Listings erlauben
	{Ö}{{\"O}}{1}
	{Ä}{{\"A}}{1}
	{Ü}{{\"U}}{1}
	{ß}{{\ss}}{1}
	%{ẞ}{{\MakeUppercase{\ss}}}{1} % KNOWNISSUE: der aktuell verwendete font enthält das zeichen leider nicht, also lieber erstmal weglassen ;-)
	{ü}{{\"u}}{1}
	{ä}{{\"a}}{1}
	{ö}{{\"o}}{1}
}

\lstdefinelanguage{thesis-latexbeispiel}
	{morekeywords=
		{\\begin, \\end, \\frac, \\pi, \\\\, \\sum, \\varphi, \\pm, \\int, \\limits, \\infty, \\cdot, \\sqrt, \\left, \\right, \\mathrm,
		\\linewidth, \\centering, \\includegraphics, \\caption, \\label, \\lstinline, \\lstinputlisting, \\todo, \\listoftodos, \\url, \\href, \\cite, \\textqoute,
		\\gls, \\glssymbol, \\glslink, \\LaTeX},
		emph={figure, table, tabular, align, lstlisting},
		sensitive=true,
		morecomment=[l]{\%},
%		morecomment=[s]{/*}{*/},
%		morestring=[b]",
		alsoletter={\\},
	}
	

%%%%%%%%%%%%%%%%%%%%%%%%%%%%%%%%%%%%%%%%%%%%%%%%%%%%%%%%%%%%%%%%%%%%%%%%%%%%%%%%
%%% Format der Float-Beschriftungen %%%%%%%%%%%%%%%%%%%%%%%%%%%%%%%%%%%%%%%%%%%%

\usepackage{caption}
\usepackage[labelformat=brace,list=true]{subcaption}
%Farbiger Text geht mit labelfont={bf,it,color=red}
% Mehr Infos zu CaptionStyle http://ctan.math.washington.edu/tex-archive/macros/latex/contrib/caption/caption-eng.pdf
\DeclareCaptionStyle{myCaptionStyle}{%
	format=hang,			% Basis-Format
	labelsep=quad,			% Abstand Typ/Nummer <-> Beschriftung
%	indention=-7em,			%
	margin=0.5em,				% Zusatzabstand Links
	singlelinecheck=false,	%
	labelfont={sf,small,bf,color=captionLabel},% Font Typ/Nummer
	textfont={sf,small, color=captionText},	% Font Beschriftung
%	justification=centering,
}

\DeclareCaptionStyle{mySubCaptionStyle}{%
	format=plain,			% Basis-Format
	labelsep=space,			% Abstand Typ/Nummer <-> Beschriftung
%	indention=-7em,			%
	margin=0.5em,			% Zusatzabstand links
	singlelinecheck=false,	%
	labelfont={sf,small,bf, color=subcaptionLabel},% Font Typ/Nummer
	textfont={sf,small, color=subcaptionText},	% Font Beschriftung
%	justification=centering,
}

% labelformat: empty  :
%              simple : a 
%              brace  : a)
%              parens : (a)


\captionsetup[table]{style=myCaptionStyle}
\captionsetup[figure]{style=myCaptionStyle}
\captionsetup[lstlisting]{style=myCaptionStyle}

\captionsetup[subfigure]{style=mySubCaptionStyle}
\captionsetup[subtable]{style=mySubCaptionStyle}


%%%%%%%%%%%%%%%%%%%%%%%%%%%%%%%%%%%%%%%%%%%%%%%%%%%%%%%%%%%%%%%%%%%%%%%%%%%%%%%%
%%% PDF-Meta-Informationen und Links im Dokument %%%%%%%%%%%%%%%%%%%%%%%%%%%%%%%
\RequirePackage{nameref}% Workaround: Loading nameref here fixes problem with package-order of nameref and showkeys

\RequirePackage[]{hyperref}

\ifbool{linksMarkieren}{
	% nichts
}{
	\hypersetup{hidelinks} % Link-Markierung verstecken
}

\hypersetup{
	linktoc=all, 						% Welcher Teil im Inhaltsverzeichnis ist klickbar? (none/section/page/all)
    colorlinks=false,					% false: boxed links (wird nicht mitgedruckt); true: colored links (wird auch so gedruckt!)
	linkcolor=primary!50!white,			% color of internal links
    linkbordercolor=primary!20!white,
%
    citecolor=primary!50!black,			% color of links to bibliography
    citebordercolor=primary!20!gray,
%
    filecolor=AccentStrong,				% color of file links
    filebordercolor=AccentStrong!20,
%
    urlcolor=AccentWeak,				% color of external links
    urlbordercolor=AccentWeak!20,
%
	menucolor=blue,
	menubordercolor=blue,
%
	runcolor=red,
	runbordercolor=red,
%
    %bookmarks=true,			% show bookmarks bar?
    unicode=true,				% non-Latin characters in Acrobat’s bookmarks
    pdftoolbar=true,			% show Acrobat’s toolbar?
    pdfmenubar=true,			% show Acrobat’s menu?
    pdffitwindow=false,			% window fit to page when opened
    pdfstartview={FitH},		% fits the width of the page to the window
    pdftitle=\thema,			% Titel
    pdfauthor=\autor,			% Autor
    pdfsubject={\artderarbeit\ am TMDT},	% subject of the document
    pdfcreator=LaTeX,			% creator of the document
    pdfproducer=,				% producer of the document
    pdfkeywords=\schlagwoerter,	% Schlagworte
    pdfnewwindow=true,			% Links in neuem Fenster öffnen
}


%%%%%%%%%%%%%%%%%%%%%%%%%%%%%%%%%%%%%%%%%%%%%%%%%%%%%%%%%%%%%%%%%%%%%%%%%%%%%%%%
%%% Glossar, Abkürzungen, Akronyme, Symbole, ... %%%%%%%%%%%%%%%%%%%%%%%%%%%%%%%

\usepackage[
	toc=false,		% Nicht im Inhaltsverzeichnis anzeigen. Wenn die Verzeichnisse dort erscheinen sollen, machen wir das selber
	translate=babel,% Babel zum Übersetzen verwenden
	acronym,		% Eigenes Verzeichnis für Akronyme	(CRC, BAföG, LASER, ...)
	abbreviations,	% Eigenes Verzeichnis für Abkürzungen (Dr., z.B., bspw., bzw., etc.)
	symbols,		% Eigenes Verzeichnis für Symbole (µ, ...)
]{glossaries-extra}

\usepackage{glossary-longragged}% Extra Style
\usepackage{glossary-mcols}% Extra Style

% 'seename' wird mit translate=babel nicht übersetzt => dann machen wir das eben selber
\DeclareTranslation{German}{seename}{siehe}
\DeclareTranslation{English}{seename}{see}
\DeclareTranslationFallback{seename}{\seename} 
\renewcommand{\seename}{\expandonce{\GetTranslation{seename}}} %expandonce: es scheint so als würde der Befehl zu früh/spät expandiert - mit expandonce passt es

\makenoidxglossaries %langsames Sortieren, funktioniert aber immer und braucht keine externen Programme
\loadglsentries{Verzeichnisse/Glossar.tex}

%%%%%%%%%%%%%%%%%%%%%%%%%%%%%%%%%%%%%%%%%%%%%%%%%%%%%%%%%%%%%%%%%%%%%%%%%%%%%%%%
%%%%%%%%%%%%%%%%%%%%%%%%%%%%%%%%%%%%%%%%%%%%%%%%%%%%%%%%%%%%%%%%%%%%%%%%%%%%%%%%
% Anpassung von Nummerierungstiefe etc. (default passt aber schon)
\setcounter{secnumdepth}{\kapitelTiefeNummerierung} %Nummerierungstiefe
\setcounter{tocdepth}{\kapitelTiefeInhaltsverzeichnis} %Kapiteltiefe im Inhaltverzeichnis



\usepackage{blindtext}

%%%%%%%%%%%%%%%%%%%%%%%%%%%%%%%%%%%%%%%%%%%%%%%%%%%%%%%%%%%%%%%%%%%%%%%%%%%%%%%%
%%% Todos %%%%%%%%%%%%%%%%%%%%%%%%%%%%%%%%%%%%%%%%%%%%%%%%%%%%%%%%%%%%%%%%%%%%%%

% Hack: todonotes mag keine \marginparwidth < 2cm => Warnung
%       In unserem Fall ist das wegen der kleinen Schrift aber OK.
%       => Vor dem Laden einen größeren Wert setzen, danach den richtigen Wert
\AtBeginDocument{%
	\setlength{\marginparwidth}{10cm}% Dummy-Wert >= 2cm
}
\usepackage[
	\hauptsprache,
%	textwidth=\randRechtsAussen-0em,%TODO sollte sich berechnen lassen, oder?
	textsize=scriptsize,
	backgroundcolor=TmdtLightBlue!20!white,
	linecolor=TmdtOrange,
	bordercolor=TmdtDarkBlue,
	textcolor=black
]{todonotes}

\AtBeginDocument{% korrekte Breite laden
	\setlength{\marginparwidth}{\randAussen-2.5mm}% hieraus ergibt sich die Breite der Todo-Boxen
}

%%%%%%%%%%%%%%%%%%%%%%%%%%%%%%%%%%%%%%%%%%%%%%%%%%%%%%%%%%%%%%%%%%%%%%%%%%%%%%%%
%%% 
\usepackage[titles]{tocloft}

\usepackage{xpatch}

% List of Listings: Indent vermeiden
\makeatletter
\xpatchcmd\l@lstlisting{1.5em}{0em}{}{} % TODO: wirkt nicht sehr robust
\makeatother

%%%%%%%%%%%%%%%%%%%%%%%%%%%%%%%%%%%%%%%%%%%%%%%%%%%%%%%%%%%%%%%%%%%%%%%%%%%%%%%%
%% Für Debugszwecke: Mehr Infos zu Boxen ausgeben:
%\showboxdepth=\maxdimen
%\showboxbreadth=\maxdimen
